%! Graphing Exponential Functions and Transformations — Unit 6, Lesson 6.2 — Lesson Plan
\documentclass[10pt]{article}
\usepackage{algebra2-article}
\usepackage{algebra2-boxes}
\usepackage{graphicx}   % -article does NOT load graphicx; needed for thumbnails/screenshots
\graphicspath{{images/}}
\newcommand{\TallMath}[1]{$\displaystyle #1\rule[-1.4em]{0pt}{3.2em}$}
\newcommand{\CourseName}{Algebra 2: Shepherd}
\newcommand{\SchoolYear}{2026--2027}
\newcommand{\MeetingLength}{55 minutes}
\newcommand{\UnitNumberName}{Unit 6: Exponential Functions \quad}
\newcommand{\LessonNumberName}{Lesson 6.2: Graphing Exponential Functions and Transformations}

\begin{document}
\begin{center}
    {\Large\bfseries \CourseName: \SchoolYear \\
    {\small\bfseries \UnitNumberName \LessonNumberName}}
\end{center}

\begin{tcolorbox}[colback=forestbg, colframe=forest, boxrule=0.9pt, arc=2mm,
    left=2mm, right=2mm, top=1.5mm, bottom=1.5mm]
    \textbf{Primary Objective:} Students can graph $g(x)=a\,b^{\,x-h}+k$ by \emph{moving the parent}
    $y=b^{x}$, and can read the horizontal asymptote, the range, and the $y$-intercept straight off
    the equation without plotting a point. The lesson turns on a single asymmetry: of the four
    transformations, \emph{only} the vertical shift $f(x)+k$ moves the asymptote --- so it is the only
    one that changes the range, and the only one that can give the graph an $x$-intercept at all.
    Students also separate the two minus signs that look alike and are not: $-b^{x}$ reflects across
    the $x$-axis (moving the graph to the other side of the asymptote), while $b^{-x}$ reflects
    across the $y$-axis and turns growth into decay, because $b^{-x}=\left(\frac1b\right)^{x}$.
    Finally, students run the reading backwards --- given a graph, recover $k$ from the asymptote, $a$
    from the $y$-intercept \emph{minus $k$}, and $b$ from one more point.
    \\[2pt]
    \textbf{Standards (2023 VA SOL):} \textbf{A2.F.1b} and \textbf{A2.F.1c} are the assessed standards
    of the day --- the transformation lens ($a$, $h$, $k$, and reflections) applied to the exponential
    parent, and the effect of each parameter on the graph. \textbf{A2.F.1e} (compare and contrast
    graphs, tables, and equations of a function family) is the structural spine: every section moves
    between an equation, a pre-drawn graph, and a table of values, and Section 4 reverses the
    direction entirely. \textbf{A2.F.2h} (horizontal asymptotes; the standard names exponential
    explicitly) and \textbf{A2.F.2a} (domain and range) are what the whole lesson is measured
    against --- ``did the asymptote move, and what is the range now?'' is the question asked on every
    item. \emph{Note:} Unit 6 has no \texttt{A2.EO} or \texttt{A2.EI} standard --- the unit rests
    entirely on \texttt{A2.F} and \texttt{A2.ST.2}.
\end{tcolorbox}

\renewcommand{\arraystretch}{1.4}
\begin{skillbox}[Priority Ideas \& Skills]{goldbox}
\begin{tabularx}{\linewidth}{@{}X|X@{}}
\textbf{Priority Ideas \& Skills} & \textbf{Key Understandings (the ``why'')} \\[2pt]
\hline
\vspace{-6pt}
\begin{itemize}[leftmargin=1.1em, nosep, after=\vspace{2pt}]
  \item Read $a$, $h$, and $k$ out of $g(x)=a\,b^{\,x-h}+k$ and name the move each one makes.
  \item Give the \textbf{horizontal asymptote} of any transformed exponential: $y=k$, and nothing else can change it.
  \item Give the \textbf{range}: $(k,\infty)$ when $a>0$, $(-\infty,k)$ when $a<0$ --- always a parenthesis at $k$.
  \item Find the $y$-intercept as $a+k$, not $a$.
  \item Distinguish $-b^{x}$ (flip over the $x$-axis) from $b^{-x}$ (flip over the $y$-axis, growth $\to$ decay).
  \item Decide whether a graph has an \textbf{$x$-intercept} from the signs of $a$ and $k$ alone.
  \item Match a pre-drawn graph to its equation using the dashed asymptote first.
  \item \textbf{Write the equation} of a transformed exponential from its graph: $k$, then $a$, then $b$.
\end{itemize}
&
\vspace{-6pt}
\begin{itemize}[leftmargin=1.1em, nosep, after=\vspace{2pt}]
  \item \emph{Only $f(x)+k$ moves the asymptote --- and everything that depends on the asymptote moves with it.} One sentence; it decides every item today.
  \item The asymptote is the \emph{boundary of the range} (Lesson 6.0). Move the boundary and the range moves; leave it and the range cannot change no matter how far sideways the graph slides.
  \item A horizontal shift is invisible to the range. That is genuinely surprising, and it is why the inside/outside distinction finally has consequences.
  \item Growth versus decay is read off $b$ against $1$ --- never off $a$. A negative $a$ chooses the \emph{side} of the asymptote, not the family.
  \item ``Exponentials never cross the $x$-axis'' was a claim about $y=ab^{x}$. It is now refined, not contradicted: crossing happens exactly when $a$ and $k$ have opposite signs.
  \item Reading a graph backwards is the same three facts in reverse order --- and the order matters, because the $y$-intercept is $a+k$.
  \item In a real model the asymptote is a physical thing (room temperature), which is why $k$ is worth this much attention.
\end{itemize}
\end{tabularx}
\end{skillbox}

\begin{skillbox}[{Vocabulary, Concepts, \& Theorems}]{sky}
\renewcommand{\arraystretch}{1.3}
\begin{tabularx}{\linewidth}{@{}l X@{}}
\textbf{Transformation form} & $g(x)=a\,b^{\,x-h}+k$ --- the parent $y=b^{x}$ with all four moves attached. \\
\textbf{Vertical shift $k$} & Slides the graph up or down \emph{together with} its asymptote, which becomes $y=k$. The only asymptote-mover in the form. \\
\textbf{Horizontal shift $h$} & $b^{\,x-h}$ slides the graph $h$ right ($h$ left if negative). Asymptote and range unchanged. \\
\textbf{Vertical stretch / compression $a$} & $|a|>1$ stretches, $0<|a|<1$ compresses. Sets the $y$-intercept; never moves the asymptote. \\
\textbf{Reflection across the $x$-axis} & $a<0$. Puts the graph on the other side of the asymptote, so the range becomes $(-\infty,k)$. The family is unchanged. \\
\textbf{Reflection across the $y$-axis} & $f(-x)$. For an exponential this \emph{changes the base}: $b^{-x}=\left(\frac1b\right)^{x}$, so growth becomes decay. \\
\textbf{Horizontal asymptote $y=k$} & The boundary the range excludes (Lesson 6.0), now movable. \\
\textbf{Range from the asymptote} & $(k,\infty)$ if $a>0$; $(-\infty,k)$ if $a<0$. Parenthesis at $k$, always. \\
\textbf{$x$-intercept condition} & The graph crosses the $x$-axis exactly when $a$ and $k$ have opposite signs (and never when $k=0$). \\
\end{tabularx}
\end{skillbox}

\begin{fixedskillbox}[Activate Prior Knowledge \& Spiral Review]{forestbg}
    The Warm-Up rehearses exactly what today needs, in the order the lesson uses it. \textbf{(1)}
    Shift directions on two \emph{familiar} parents --- $y=(x-3)^{2}+1$ and $y=\sqrt{x+2}-5$ ---
    finishing with the general statement that outside moves it vertically and inside moves it
    horizontally. This is the only part of today that is not new, and it is the part students must
    bring intact, because the whole lesson is about which of those two moves the asymptote.
    \textbf{(2)} A recall check on $y=3(2)^{x}$ --- $a$, $b$, growth/decay, $y$-intercept, asymptote,
    range --- because every answer today is measured against ``$y=0$ and $(0,\infty)$,'' and a student
    who cannot produce the parent's pair cannot say what changed. \textbf{(3)} $2^{-2}$ against
    $\left(\frac12\right)^{2}$ and $2^{-3}$ against $\left(\frac12\right)^{3}$, which earns
    $2^{-x}=\left(\frac12\right)^{x}$ \emph{numerically} before Section 3 states it as the reflection
    rule. Debrief item 1 at the board and leave ``outside $=$ vertical, inside $=$ horizontal'' written
    where it will stay all period.
\end{fixedskillbox}

\begin{skillbox}[Hook]{forestbg}
    Three graphs on the board, all built from $y=2^{x}$: \textbf{I} is the parent, \textbf{II} is
    $y=2^{x}-3$, and \textbf{III} is $y=2^{\,x-3}$ --- each drawn with its horizontal asymptote as a
    dashed line. Do not label them. Ask the class to match the three equations to the three graphs,
    and then ask the question the whole lesson hangs on: \textbf{``Two of these equations contain a
    $3$, and both of those graphs moved. Why is only \emph{one} of the dashed lines somewhere new?''}
    Take answers until someone connects it to the Warm-Up: the $3$ that is \emph{outside} carries the
    whole picture --- asymptote included --- down with it, and the $3$ that is \emph{inside} only
    slides the picture sideways past an asymptote that never budged. Then close the trap: point at
    graph II crossing the $x$-axis and ask, \textbf{``In Lesson 6.0 you proved an exponential can
    never cross the $x$-axis. Explain that graph.''} Let the class resolve it --- 6.0's claim was
    about $y=ab^{x}$, and the vertical shift is exactly what changes it. Do \emph{not} generalize to
    $a\,b^{\,x-h}+k$ yet; Section 1 earns the form.
\end{skillbox}

\begin{skillbox}[Lesson]{forestbg}
\begin{multicols}{2}
\textbf{1. One form, four jobs.} Write $g(x)=a\,b^{\,x-h}+k$ and fill the four-row table together.
Three of the four columns are review from Units 2 and 5; the fourth --- \emph{does the asymptote
move?} --- is the new content, and it reads YES, no, no, no. Say the asymmetry out loud rather than
letting students infer it. Then work $g(x)=2^{\,x-1}+3$ end to end: parent, moves, asymptote $y=3$,
range $(3,\infty)$, and the $y$-intercept $g(0)=2^{-1}+3=3.5$ --- which is already the lesson's
sneakiest fact, since the intercept is neither $a$ nor $k$. Close the section by naming what the
shift right did to the range: nothing.

\textbf{2. Why $k$ takes the range with it.} Anchor on $y=2^{x}-4$ with its graph. The whole picture,
asymptote included, slid down $4$: asymptote $y=-4$, range $(-4,\infty)$, $y$-intercept $(0,-3)$, and
--- the payoff --- an \textbf{$x$-intercept at $(2,0)$}, verified by $2^{2}-4=0$. Spend real time on
the red box. The bracket error ($[k,\infty)$) is Lesson 6.0's error wearing new clothes and will
reappear on the exit ticket. The second bullet is the one to say slowly: 6.0's ``no zeros, ever'' is
being \emph{refined}, and the refined rule is that the graph crosses exactly when $a$ and $k$ have
opposite signs.

\columnbreak

\textbf{3. Two minus signs.} Put $y=-2^{x}$ and $y=2^{-x}$ side by side and fill the comparison table.
The $x$-axis flip moves the graph to the other side of an unmoved asymptote --- range $(-\infty,0)$,
but the base is still $2$, so it is \emph{still growth}, just falling. The $y$-axis flip leaves the
range alone and changes the family, because $2^{-x}=\frac{1}{2^{x}}=\left(\frac12\right)^{x}$ --- the
Warm-Up's arithmetic, now a rule. Neither one touched the asymptote, because neither one is a vertical
shift. Expect ``$-2^{x}$ is decay''; refuse it and point back at 6.0.

\textbf{4. Reading backwards.} Three steps, \emph{in order}: the dashed line gives $k$; the
$y$-intercept is $a+k$, so subtract before claiming $a$; one more point gives $b$. Work the graph with
asymptote $y=2$ through $(0,5)$ and $(1,8)$ to $y=3(2)^{x}+2$, then check it at $x=2$. Finish with the
warning: a student who skips step 1 writes $a=5$ here, every time. This is the skill Lesson 6.5 leans
on. Preview Lesson~6.3.
\end{multicols}
\end{skillbox}

\begin{skillbox}[Active Monitoring]{redbox}
Circulate during the notes practice and Tier work. Watch for and correct:
\begin{itemize}[leftmargin=1.2em, nosep]
  \item treating an \textbf{inside} number as a vertical shift --- ``$y=5^{\,x-2}$ has asymptote $y=-2$''; this is \emph{the} signature error of the lesson;
  \item claiming the range changed after a \textbf{horizontal} shift, or after a stretch;
  \item the \textbf{bracket} error --- writing the range as $[k,\infty)$; the asymptote is excluded, always;
  \item reading the $y$-intercept as $a$ instead of $a+k$ (and the reverse: writing $a$ for the intercept when going backwards);
  \item calling $-b^{x}$ \textbf{decay}; the base decides the family, the sign of $a$ decides the side of the asymptote;
  \item confusing $-b^{x}$ with $b^{-x}$ --- ask which axis, not which sign;
  \item the horizontal-shift direction ($b^{\,x-3}$ moves \emph{right}), the Unit 2 / Unit 5 sign flip;
  \item hunting for an $x$-intercept on a graph that has none, or missing one on a graph that does --- send them to the signs of $a$ and $k$;
  \item in Tier E, concluding that $2^{\,x-3}=\frac18\cdot2^{x}$ means shifts and stretches are interchangeable \emph{in general} --- part (d) exists to kill that.
\end{itemize}
Cold-call prompts: ``Did the asymptote move? Say why or why not.'' ``What is the range now --- and is
that a bracket or a parenthesis?'' ``Your $a$ and the $y$-intercept are different numbers. Which is
which?'' ``Is that a negative on the input or on the output?'' ``Does this graph cross the $x$-axis?
Answer from the signs, not the picture.''
\end{skillbox}

\begin{skillbox}[Group Work \& Differentiation]{redbox}
\textbf{Move It --- and Say What Moved} activity (tiered; mirrors the notes).
\begin{multicols}{3}
\textbf{Tier R --- Remediate.} A four-row ``what moved?'' table off the parent $y=3^{x}$ (up $4$,
right $2$, down $5$, flip over the $x$-axis) closing with the question that \emph{is} the lesson:
which two rows kept the parent's asymptote, and what do they have in common? Then one function taken
all the way through --- $g(x)=2^{x}+3$: parent, move, asymptote, a table of five values, the
$y$-intercept, the range, why no input ever gives exactly $3$, and why there is no $x$-intercept.

\columnbreak
\textbf{Tier A --- Approaching.} Four pre-drawn graphs of transformations of $y=2^{x}$ matched to
their equations, each with asymptote and range read from the \emph{equation}, and then the diagnostic
question: exactly one of the four has an $x$-intercept --- which, and why is it the only one? Next,
graph $\to$ equation with a \emph{negative} $k$ (asymptote $y=-1$, intercept $(0,2)$), where $a=3$
comes out \emph{larger} than the intercept. Then the \textbf{Justify}: Nadia's claim that $-5^{x}$ and
$5^{-x}$ are ``the same thing, both have a negative'' --- axis, range, and family required for each.

\columnbreak
\textbf{Tier E --- Extension.} Three parts. \emph{A shift in disguise:} $2^{\,x-3}=\frac18\cdot2^{x}$
proved by exponent rules and checked in a table --- a horizontal shift and a vertical compression are
the \emph{same} move for this family --- with part (d) testing the same trick on $(x-3)^{2}$, where it
fails. \emph{All four moves at once:} $g(x)=-2(3)^{\,x-1}+5$, including an exact $y$-intercept of
$\frac{13}{3}$ and an $x$-intercept decided from signs. \emph{Coffee:} $T(t)=68+132(0.85)^{t}$, where
the asymptote is \textbf{room temperature} --- with an algebraic-versus-contextual range beat and a
closing $132(0.85)^{t}=32$ that has no common base.
\end{multicols}
\end{skillbox}

\begin{skillbox}[Individual Work \& Assessment]{redbox}
\textbf{Exit Ticket} (independent, no notes), four items on one page: \textbf{(1)} read $y=4^{x}-5$
for asymptote, range, $y$-intercept, and family --- the family answer is \emph{growth}, a deliberate
check that students are still reading $b$ and not the negative intercept; \textbf{(2)} the
inside/outside discrimination --- which of $y=2^{\,x+3}$ and $y=2^{x}+3$ has range $(3,\infty)$, plus
one sentence on what the other does instead; \textbf{(3)} write the equation from a pre-drawn graph
(asymptote $y=1$, $y$-intercept $(0,3)$, through $(1,5)$), the hardest and most informative item, where
$a=3$ is the error to look for; and \textbf{(4)} an \textbf{SOL-style multiple-choice} item --- the
range of $g(x)=3(2)^{x}-6$ --- whose distractors are the three errors of the day: $[-6,\infty)$ is the
bracket error, $(0,\infty)$ ignored the shift, and $(6,\infty)$ lost the sign of $k$. Sort the collected
tickets into \emph{got it / inside-outside / bracket / grabbed the intercept}. Only the inside-outside
pile needs a full re-teach, and it should be run off the Hook's graphs II and III, where the same $3$
does two different things.
\end{skillbox}

\begin{skillbox}[Reinforcement \& Extension]{goldbox}
\textbf{Homework:} a five-row feature table (asymptote / range / $y$-intercept / family) whose fourth
row, $y=-4^{x}$, is the ``still growth, but falling'' trap, closing with \emph{which rows have an
$x$-intercept?} --- answer: only the one where $a$ and $k$ have opposite signs; five transformations
described in words with their new asymptotes; a four-graph matching set where \emph{two} graphs cross
the $x$-axis and both cross at $x=2$ (deliberately, since both reduce to $2^{x}=4$); two graph-to-equation
problems, the second with $b=\frac12$ coming out of $6b+1=4$, which students will not expect; and an
\textbf{error-analysis} pair carrying the lesson --- Marcus reading $y=5^{\,x-2}$ as a vertical shift,
and Tara calling both $-6^{x}$ and $6^{-x}$ decay. \textbf{Extension:} a soda at $38^{\circ}$F warming
in a $72^{\circ}$ room, $T(t)=72-34(0.9)^{t}$ --- $a$ is \emph{negative} because the soda starts below
the asymptote, the asymptote is room temperature, the algebraic range $(-\infty,72)$ differs from the
contextual $[38,72)$, and the $x$-intercept genuinely exists but sits at negative $t$, which students
argue \emph{without} logarithms (a base under $1$ exceeds $1$ only at a negative exponent). Do not name
logarithms; that is Unit 7. \textbf{Preview:} Lesson~6.3, \emph{Solving Exponential Equations with a
Common Base} --- where the $x$-intercept students found by \emph{looking} at $y=2^{x}-4$ gets solved,
because $2^{x}=4$ has a common base. The coffee and soda equations do not, and both are left standing.
\end{skillbox}
\end{document}
